\bilingualchapter{Asset Allocating Investor}{资产配置型投资者}

\begin{englishblock} THIS CHAPTER IS FOR ASSET ALLOCATING INVESTORS WHO MOSTLY don't believe that assets' prices can be forecasted, and use the ‘no-rule' trading rule within a systematic framework to allocate trading capital between different assets. \end{englishblock}

\begin{chineseblock}本章写给资产配置型投资者，他们大体上并不相信资产价格是可以被预测的，而是使用“无规则”交易规则，并把它放进系统化框架中，以便在不同资产之间分配交易资本。\end{chineseblock}

\bilingualsection{Chapter overview}{本章概览}

\begin{bilingualtableblock} \begin{tabularx}{\textwidth}{@{} L{0.30\textwidth} Y@{}}\textbf{Who are you?} & Introducing the asset allocating investor. \\

\textbf{你是谁？} & 介绍“资产配置型投资者”这一类人。 \\
\textbf{Using the framework} & How you will use the systematic framework for asset allocation. \\
\textbf{使用框架} & 如何使用系统化框架来做资产配置。 \\
\textbf{Weekly process} & The weekly process that asset allocators should use. \\
\textbf{每周流程} & 资产配置者应当采用的每周流程。 \\
\textbf{Trading diary} & A diary showing how you could have invested over a few hectic weeks of 2008. \\
\textbf{投资日记} & 一份日记，展示你在 2008 年那几周剧烈波动时期里 可能会如何进行投资。 \\
\end{tabularx}
\end{bilingualtableblock}

\bilingualsection{Who are you?}{你是谁？}

\begin{englishblock} If you're an asset allocating investor you believe the best returns can be obtained by investing in a diversified portfolio of assets without trying to predict relative risk adjusted returns. The systematic framework in this book will allow you to do that, using the ‘no rule' rule I introduced in chapter seven, ‘Forecasts'. This rule has a constant forecast of +10, implying you think all assets will have the same Sharpe ratio. To get access to different asset classes I assume you'll be using exchange traded funds (ETFs), which are relatively cheap collective funds. These will allow you to access a wide range of asset classes whilst only holding positions in a relatively small number of \end{englishblock}

\begin{chineseblock}如果你是一名资产配置型投资者，那么你相信：最好的收益，来自于持有一个充分分散的资产组合，而不是去尝试预测各类资产之间风险调整后收益的相对高低。本书中的系统化框架就能帮助你做到这一点，你会用到我在第七章“预测值”中介绍的那条“无规则”规则。这条规则使用一个恒定的预测值 +10，也就是说，你在默认假设所有资产都拥有相同的夏普比率。为了接触不同资产类别，我这里假设你使用的是交易所交易基金（ETF），这是一种相对便宜的集合基金。它们能让你在只持有相对较少数量品种的前提下，覆盖广泛的资产类别。\end{chineseblock}

\bipairgap

\begin{englishblock} instruments. I am assuming that you're not comfortable, or unable, to use leverage or to do short selling.148 In the example of this chapter I use a notional trading capital of €10,000,000, the size of a small pension fund, although everything in this chapter is relevant to those with much larger and smaller portfolios. Your imaginary pension fund trustees have specified some constraints on the portfolio. You can only invest in bonds and equities, and no more than 40\% of your capital can be allocated to bonds. Within the equity portfolio no more than 30\% can be put into emerging markets. In bonds you can't put more than 25\% into emerging markets and no more than 25\% into inflation linked bonds. All constraints are expressed in risk adjusted proportions. An extension to the basic example will be to adjust instrument weights according to some expectations about asset Sharpe ratios. These aren't produced by a systematic trading rule, but this setup is common in institutional investing when allocations need to be adjusted to reflect in-house or consultants' strategic views. \end{englishblock}

\begin{chineseblock}这里我假设，你既不愿意、或者也无法使用杠杆，也不会进行卖空。148 在本章示例中，我使用一个名义上的 €10,000,000 交易资本，这大致相当于一个小型养老金基金的规模；不过，本章讨论的所有内容，其实也适用于规模比这大得多或小得多的投资组合。你那虚构的养老金基金受托人，对投资组合设置了一些约束条件：你只能投资债券和股票，而且用于债券的资本配置不得超过总资本的 40\%。在股票部分中，新兴市场的配置不得超过 30\%。在债券部分中，新兴市场债券的配置不得超过 25\%，通胀挂钩债券的配置也不得超过 25\%。所有这些约束，都是以风险调整后的比例来表达的。这个基础示例还可以进一步扩展：你可以根据对资产夏普比率的某些预期，来调整品种权重。这些预期并不是由系统化交易规则生成的，但在机构投资中，这种设置非常常见，因为投资组合配置经常需要反映内部团队或外部顾问所给出的战略观点。\end{chineseblock}

\bilingualsection{Using the framework}{使用框架}

\begin{englishblock} Instrument choice: diversification, costs, volatility and size You have quite a large account size, so in theory you could have an allocation to a large number of ETFs without having any difficulties with the minimum size problem I discussed in chapter twelve, ‘Speed and Size' (page 200). However to make this example more tractable I am going to limit your portfolio to just ten instruments. In selecting ETFs I've generally gone for those with the lowest total expense ratio (this is an annual holding cost you pay regardless of how often you trade). Additionally, because you can't use leverage there would be issues hitting the desired percentage volatility target if your portfolio contains very low volatility assets, as you'll see below. In particular you should be concerned about bond ETFs that could have a sigma below 5\% per year. To avoid this problem I've generally suggested longer maturity bond funds on which the volatility will be higher. Low risk instruments have other problems of course, which I've discussed earlier in the book, including generally higher standardised costs. I'll assume, as I did in chapter twelve, ‘Speed and Size', that you trade in 100 share blocks as this is more economical. Using this size you could implement this strategy with just €100,000 of trading capital and no issues of minimum size.149 \end{englishblock}

\begin{chineseblock}品种选择：分散化、成本、波动率与规模 你的账户规模相当大，因此理论上你完全可以把资金分配到许多 ETF 上，而不会遇到第十二章“速度与规模”（第 200 页）中讨论过的最小交易单位问题。不过，为了让这个例子更容易处理，我会把投资组合限制在仅 10 个品种之内。在选择 ETF 时，我通常优先选择总费用率最低的产品（这是你每年都要支付的持有成本，无论你交易得多频繁）。此外，由于你不能使用杠杆，如果投资组合中包含波动率极低的资产，那么你就会在实现目标百分比波动率方面遇到困难，正如你很快会看到的那样。尤其需要警惕的是，那些年化 sigma 可能低于 5\% 的债券 ETF。为了避免这个问题，我通常会建议使用久期更长的债券基金，因为它们的波动率会更高。当然，低风险工具还有其他问题，我在本书前面已经讨论过了，其中包括通常更高的标准化成本。和第十二章“速度与规模”中一样，我在这里也假设你按 100 股一个单位来交易，因为这样更经济。在这个单位下，你其实只用 €100,000 的交易资本就能实施这个策略，而且不会遇到最小交易单位问题。149 \end{chineseblock}

\begin{enumerate}
\item This doesn't preclude buying short ETFs or those with leverage baked in, although I won't be using them in this example. \end{enumerate}

\begin{enumerate}
\item 这并不意味着你绝对不能买入做空型 ETF 或带有内置杠杆的 ETF，只是我在本例中不会使用它们。\end{enumerate}

\begin{englishblock} The costs of ETFs will vary depending on their volatility, so we'll conservatively use the standardised cost for the lowest risk ETF (IGIL: global inflation linked bonds) for all instruments. The relevant cost of 0.08 Sharpe ratio units was calculated in chapter twelve, ‘Speed and Size'. There are several other considerations when choosing ETFs and you should consult one of the books mentioned in appendix A to understand these. Finally, this selection does not constitute an endorsement for these specific products, and the information I've used to choose them will almost certainly change in the future. Table 39 shows the final set of instruments. In the example you're constrained in your portfolio allocations by an outside investor. You have limits on how much you can allocate to bonds and equities, and to emerging markets and inflation linked bonds. This is reflected in the asset grouping I've chosen in the table. You need the grouping to allocate the instrument weights in your portfolio, which I'll do below using the handcrafting method. \end{englishblock}

\begin{chineseblock} ETF 的成本会随着其波动率不同而变化，因此我们在这里采取保守做法：对所有品种都使用风险最低的那只 ETF （IGIL：全球通胀挂钩债券）对应的标准化成本。这个 0.08 个夏普比率单位的成本值，是在第十二章“速度与规模”中算出来的。在选择 ETF 时，还有其他若干考虑因素，你应当去看附录 A 中提到的相关书籍，以更好地理解这些问题。最后，这里给出的这些具体产品并不构成我对它们的推荐，而且我用来选择它们的信息在未来也几乎肯定会发生变化。表 39 展示了最终选择的品种集合。在这个例子中，你的投资组合配置受到外部投资者的约束。你在债券、股票，以及新兴市场和通胀挂钩债券上的配置比例，都受到上限限制。这些限制被直接体现在表格中我所选定的资产分组方式上。而你正需要这些分组，才能进一步用手工构造法来配置投资组合中的品种权重，下面我就会这么做。\end{chineseblock}

\begin{bilingualtableblock} \rawtablecaption{WHAT ETFS ARE YOU TRADING AND HOW SHOULD THEY BE GROUPED?\\你在交易哪些 ETF？它们又应该如何分组？}\noindent{\small \setlength{\tabcolsep}{4pt} \renewcommand{\arraystretch}{1.12} \begin{tabularx}{\textwidth}{@{} p{0.16\textwidth} p{0.20\textwidth} Y@{}}
\toprule
\textbf{Asset class} & \textbf{Region/sub-class} & \textbf{ETFs (Ticker)} \\
\textbf{资产类别} & \textbf{区域/子类别} & \textbf{ETF（代码）} \\
\midrule
Bonds & Developed & \$ US Bonds (IDTL) \\
债券 & 发达市场 & \$ 美国债券 (IDTL) \\
Bonds & Developed & € Euro bonds (EXHK) \\
债券 & 发达市场 & € 欧元区债券 (EXHK) \\
Bonds & Developed & £ UK bonds (VGOV) \\
债券 & 发达市场 & £ 英国债券 (VGOV) \\
Bonds & Emerging & \$ EM government bonds (SEMB) \\
债券 & 新兴市场 & \$ 新兴市场政府债券 (SEMB) \\
Bonds & Inflation linked & \$ Global inflation (IGIL) \\
债券 & 通胀挂钩 & \$ 全球通胀挂钩债券 (IGIL) \\
\addlinespace[0.4em]
Equities & Developed & € Euro Stoxx 50 (EXW1) \\
股票 & 发达市场 & € Euro Stoxx 50 (EXW1) \\
Equities & Developed & \$ S\&P 500 ETF (VUSD) \\
股票 & 发达市场 & \$ S\&P 500 ETF (VUSD) \\
Equities & Developed & £ UK equities (ISF) \\
股票 & 发达市场 & £ 英国股票 (ISF) \\
Equities & Developed & € Japan equities, € hedged (EXX7) \\
股票 & 发达市场 & € 日本股票，欧元对冲 (EXX7) \\
Equities & Emerging & \$ MSCI EM ETF (VDEM) \\
股票 & 新兴市场 & \$ MSCI 新兴市场 ETF (VDEM) \\
\bottomrule
\end{tabularx}
}
\end{bilingualtableblock}

\begin{englishblock} The table shows the structure of the Asset Allocating Investor example portfolio. The grouping is driven by allocation constraints and the correlations in table 40. \end{englishblock}

\begin{chineseblock}这张表展示的是资产配置型投资者示例组合的结构。分组方式由配置约束和表 40 中的相关性共同决定。\end{chineseblock}

\bilingualsubsection{The no-rule trading rule, position sizing and costs}{无规则交易规则、仓位规模设定与成本}

\begin{englishblock} Because you're going to use the ‘no-rule' trading rule I introduced in chapter seven your forecast will always be +10. To scale positions you'll need to estimate the standard deviation of instrument returns, the price volatility. As suggested in chapter ten, ‘Position Sizing', you'll estimate this with a moving average of daily returns. However you need to think about what look-back to use for the moving average. Because ETFs are quite expensive you should follow the advice in chapter twelve, ‘Speed and Size' (page 196), and use a relatively slow 20 week (100 business day) look-back to estimate price volatility.\footnote{Alternatively using an exponentially weighted moving average estimate with a look-back of 144 days is equivalent to a 100 day simple moving average. \par 另一种做法，是使用回看期为 144 天的指数加权移动平均估计，它大致等价于 100 天的简单移动平均。} In the earlier chapter I calculated that assuming conservative costs for trading ETFs the longer look-back would cost you 0.032 Sharpe ratio (SR) units annually, compared to 0.13 SR units for the default look-back of five weeks, which is a considerable saving. With a 20 week look-back the turnover of your trades, in round trips per year, will be 0.4 times a year. If you use the speed limit I recommended in chapter twelve, and do not spend more than 0.08 SR each year on costs, then a turnover of 0.4 round trips per year implies any instrument you trade needs to have a maximum standardised cost of 0.08 ÷ 0.4 = 0.20 SR units per year. Fortunately by avoiding low volatility bonds all of your instruments are significantly cheaper than this, with the inflation linked bond ETF IGIL the most expensive with a standardised cost of 0.08 SR units. \end{englishblock}

\begin{chineseblock}由于你会使用我在第七章中介绍的“无规则”规则，因此你的预测值始终都是 +10。为了给仓位定规模，你需要估计品种收益的标准差，也就是价格波动率。按照第十章“仓位规模设定”中的建议，你可以用日收益的移动平均来估计它。不过，你还需要考虑这个移动平均应该使用多长的回看窗口。因为 ETF 的交易成本相对较高，所以你应当遵循第十二章“速度与规模”（第 196 页）中的建议，使用较慢的 20 周（100 个交易日）回看窗口来估计价格波动率。在前面的章节中我算过：如果对 ETF 交易采用保守成本假设，那么较长回看窗口对应的年度成本大约是 0.032 个夏普比率（SR）单位，而默认五周窗口则是 0.13 个 SR 单位，这节省非常可观。采用 20 周回看窗口时，你的交易换手率大约会是每年 0.4 次往返交易。如果你采用第十二章中我建议的速度上限，也就是每年在成本上不超过 0.08 个 SR，那么每年 0.4 次往返交易就意味着你所交易的任何工具，其最大标准化成本必须不高于0.08 ÷ 0.4 = 0.20 个 SR 单位每年。幸运的是，只要你避开了低波动债券，那么你所有工具的成本都会明显低于这个阈值；其中最贵的是通胀挂钩债券 ETF IGIL，它的标准化成本是 0.08 个 SR 单位。\end{chineseblock}

\bilingualsubsection{Volatility target calculation}{波动率目标的计算}

\begin{englishblock} I've set your trading capital at €10,000,000, but how do you follow the advice in chapter nine to set your percentage volatility target? In this example it's most likely to be constrained by your lack of leverage, rather than a tolerance for losses or Sharpe ratio expectations. Working out the achievable volatility target given limited leverage is a two step process. The first step involves making an initial guess as to what the target should be. You then work out if that is higher or lower than you could achieve given the amount of leverage you have available. The initial guess is then adjusted. \end{englishblock}

\begin{chineseblock}我把你的交易资本设定为 €10,000,000，但你要如何按照第九章中的建议来设定百分比波动率目标呢？在这个例子里，最可能约束你的并不是对亏损的承受能力，也不是你对夏普比率的预期，而是你根本不使用杠杆。要在杠杆受限的情况下算出可实现的波动率目标，需要两步。第一步，是先对目标值做一个初始猜测。然后，你再去判断：在当前能使用的杠杆水平下，这个猜测是太高了，还是太低了。最后，再据此调整这个初始猜测。\end{chineseblock}

\bipairgap

\begin{englishblock} Desired leverage I define the leverage factor as the size of your portfolio, divided by factor the value of your trading capital. So the factor would be 50\% if you've invested half your trading capital and kept the rest in cash, 100\% if you are using all your cash with no leverage, 200\% if you've borrowed to invest twice your trading capital, and so on. First work out what your desired leverage factor would be. You should set this just below the maximum you can achieve. It should be less than the maximum because if price volatility falls you'll want to buy more assets, so it's worth keeping some cash in reserve. I recommend a desired leverage factor of 90\% for investors who don't use leverage. \end{englishblock}

\begin{chineseblock}目标杠杆因子 我把杠杆因子定义为：投资组合总规模除以交易资本价值。因此，如果你只投资了一半资本、另一半保留为现金，那么杠杆因子就是 50\%；如果你不用杠杆、把全部现金都投进去，那么就是 100\%；如果你借钱把投资规模做到资本的两倍，那么杠杆因子就是 200\%，以此类推。第一步，是先想清楚你希望的目标杠杆因子是多少。它应当略低于你实际能够达到的最大值。之所以要低一点，是因为如果价格波动率下跌，你会希望买入更多资产，因此预留一些现金缓冲是有价值的。对于不用杠杆的投资者，我建议把目标杠杆因子设为 90\%。 \end{chineseblock}

\bipairgap

\begin{englishblock} Annualised For each instrument multiply the daily percentage price volatility by 16 percentage to get an annualised version. (As usual to annualise a daily volatility you volatility multiply by the ‘square root of time', which with around 256 business days in a year is 16.) Annualised price volatility for the ETFs I've selected ranges from 6.88\% for inflation linked bonds, up to 22.4\% for emerging market equities. \end{englishblock}

\begin{chineseblock}年化百分比波动率 对每个品种，把日百分比价格波动率乘以 16，就能得到其年化版本。（像往常一样，把日波动率年化时需要乘以“时间平方根”；如果一年大约有 256 个交易日，那么这个数就是 16。）在我所选用的 ETF 中，年化价格波动率从通胀挂钩债券的 6.88\% 一直到新兴市场股票的 22.4\% 不等。\end{chineseblock}

\bipairgap

\begin{englishblock} Initial guess Your initial guess of the correct percentage volatility target is equal to the lowest annualised price volatility of any instrument. This is 6.88\% for the IGIL inflation linked bond. \end{englishblock}

\begin{chineseblock}初始猜测 你对正确百分比波动率目标的初始猜测，应当等于所有品种中最低的年化价格波动率。在这个例子里，这个值是通胀挂钩债券 ETF IGIL 的 6.88\%。 \end{chineseblock}

\bipairgap

\begin{englishblock} Work out starting You should now run the usual calculations from chapters ten and eleven: positions first work out the volatility scalar, and then given each instrument has a fixed forecast of 10 you can calculate the subsystem position. You also use the instrument weights and the instrument diversification multiplier that are calculated later in the chapter, to get the rounded portfolio instrument position for each instrument. \end{englishblock}

\begin{chineseblock}算出初始仓位 现在你应当像第十章和第十一章中那样，完成一遍标准计算流程：先算出波动率缩放因子，然后因为每个品种都使用固定预测值 10，于是你就能得到各自的子系统仓位。接着，你还需要使用本章稍后会算出来的品种权重和品种分散化乘数，从而得到每个品种四舍五入后的投资组合仓位。\end{chineseblock}

\bipairgap

\begin{englishblock} Value of each For each instrument you need to work out the value of the starting position position. This is equal to the rounded portfolio instrument position in blocks, multiplied by 100 times the block value (since each block value represents 1\% of the value of one block of the instrument), and the usual exchange rate you use to calculate the volatility scalar. For example with your initial guess of a 6.88\% volatility target the portfolio instrument position for the UK equity ETF is 721 blocks, each of 100 shares. The value is: position 721 multiplied by block value £7.00, by 100, and the exchange rate (GBP/EUR) 1.2 is €605,640. \end{englishblock}

\begin{chineseblock}每个仓位的价值 对每个品种，你都需要算出初始仓位的实际价值。这个值等于四舍五入后的投资组合品种仓位（以 block 为单位），再乘以 100 倍的单位价值（因为每个单位价值本身表示的是该品种一个 block 价格变动 1\% 对应的金额），最后再乘以你平时用来计算波动率缩放因子的汇率。例如，若你的初始猜测波动率目标是 6.88\%，那么英国股票 ETF 的投资组合仓位会是 721 个 block，每个 block 为 100 股。对应的价值就是：721 乘以单位价值 £7.00，再乘以 100，最后再乘以汇率（GBP/EUR）1.2，得到 €605,640。 \end{chineseblock}

\bipairgap

\begin{englishblock} Total value You add up the value of all your initial positions. In this case it comes in at around €8,262,000. \end{englishblock}

\begin{chineseblock}总价值 把所有初始仓位的价值加总起来。在本例中，总和大约是 €8,262,000。 \end{chineseblock}

\bipairgap

\begin{englishblock} Calculate realised Your realised leverage will be the total value of your positions, divided leverage factor by the trading capital. Here it's €8,262,00 divided by €10,000,000 equals 0.826, or 82.6\%. This is less than one, indicating you aren't using any leverage and have some excess cash. \end{englishblock}

\begin{chineseblock}计算实际杠杆因子 你的实际杠杆因子等于总持仓价值除以交易资本。在这里，它就是 €8,262,000 ÷ €10,000,000 = 0.826，也就是 82.6\%。这小于 1，说明你没有使用任何杠杆，同时手里还留有一些多余现金。\end{chineseblock}

\bipairgap

\begin{englishblock} Adjust You can now adjust your initial guess of percentage volatility target percentage according to whether your guess is overshooting or undershooting the volatility target leverage that you want to achieve. for desired In this case you can increase your volatility target. Your desired leverage leverage was 90\%, but with the initial guess you got 82.6\%. The initial volatility target can be multiplied by 90\% ÷ 82.6\% = 1.089 So the achievable percentage volatility target is equal to the initial guess 6.88\%, multiplied by 1.089, equals 7.5\%. \end{englishblock}

\begin{chineseblock}调整百分比波动率目标 现在，你可以根据初始猜测相对于你想实现的杠杆目标是偏高还是偏低，来调整这个波动率目标。在本例中，你可以把波动率目标往上调。你的目标杠杆是 90\%，但根据初始猜测得到的实际杠杆只有 82.6\%。因此，初始波动率目标可以乘以90\% ÷ 82.6\% = 1.089。于是，可实现的百分比波动率目标就等于初始猜测 6.88\% 再乘以 1.089，最终得到 7.5\%。 \end{chineseblock}

\bipairgap

\begin{englishblock} This implies your annualised cash volatility target should be 7.5\% multiplied by your trading capital, or €750,000. This shouldn't present problems to anybody's risk appetite; and if you assume the volatility target is at the correct Half-Kelly level then the required Sharpe ratio (SR) of 0.15 should be easily achievable. It also implies you should have very modest expectations for account growth. For asset allocating investors I'd conservatively expect a maximum after cost SR of 0.4 (see page 46) which equates to returns of just 0.4 × 7.5\% = 3.0\% a year, plus the risk free interest rate (around 0.5\% in Europe as I write this), which you should include as you're not using derivatives. To improve this you'd either have to use leverage, or reduce diversification and Sharpe ratio by lowering the instrument weight on low volatility instruments like bonds (or exclude them entirely).\footnote{You can also use ETFs with internal leverage to help solve this problem, although you should first understand the complexities of their daily regearing. \par 你也可以使用带有内部杠杆的 ETF 来帮助缓解这个问题，不过在此之前，你应当先理解它们每日重新加杠杆所涉及的复杂性。} How much would you be paying in costs? I worked out above that you should expect to have a turnover of 0.4 round trips per year and standardised costs of 0.08 SR, implying an expected cost of 0.4 × 0.08 = 0.032 SR annually. With 7.5\% annualised volatility that equates to a modest annual performance drag of 0.032 × 7.5\% = 0.24\%. In practice you should pay less since I used the most expensive instrument to work out the standardised cost. On top of this would be the annual holding costs for ETFs -- currently these range from about 0.05\% for the cheapest S\&P 500 tracker up to 0.55\% for emerging market equities. \end{englishblock}

\begin{chineseblock}这意味着，你的年化现金波动率目标应当是 7.5\% 乘以交易资本，也就是 €750,000。这个水平应该不会对任何人的风险偏好构成问题；如果你进一步假设这个波动率目标正好位于正确的半 Kelly 水平，那么所要求的夏普比率（SR）仅仅是 0.15，这应该很容易实现。与此同时，这也意味着你对账户增长速度的预期应当非常克制。对于资产配置型投资者，我会保守地认为成本后的最大 SR 为 0.4（见第 46 页），这对应的收益仅仅只有 0.4 × 7.5\% = 3.0\% 每年，再加上无风险利率（在我写作时欧洲大约是 0.5\%），而由于你并没有使用衍生品，这部分无风险利率应当计入收益。如果你想把收益水平再提高，那么要么你得使用杠杆，要么你就得减少分散化，并通过降低低波动品种（比如债券）的权重（甚至彻底剔除它们）来提高夏普比率。这一结构下的成本会是多少呢？前面我算过：你的换手率大约是每年 0.4 次往返交易，标准化成本是 0.08 SR，因此预期年度成本就是 0.4 × 0.08 = 0.032 SR。以 7.5\% 的年化波动率来看，这对应的年度表现拖累大约是 0.032 × 7.5\% = 0.24\%，属于相当温和的水平。实际上，你支付的成本还应当更低一些，因为我在计算标准化成本时，使用的是最贵的那个工具。除此之外，还要再加上 ETF 本身的年度持有成本--- 目前，从最便宜的 S\&P 500 跟踪 ETF 的大约 0.05\%，一直到新兴市场股票 ETF 的大约 0.55\% 不等。\end{chineseblock}

\bilingualsubsection{Portfolios: The basics}{投资组合：基础部分}

\begin{englishblock} As I discussed in chapter eleven, a trading system consists of a portfolio of trading subsystems, one per instrument. You need to determine the instrument weights to allocate across this portfolio. Because you're using a static trading rule determining these weights is a key decision. I introduced two methods for portfolio allocation in chapter four. Bootstrapping would be an excellent way of doing this, but I'll stick to the simpler handcrafting method in this example, using rule of thumb rather than back-tested correlations. \end{englishblock}

\begin{chineseblock}正如我在第十一章中讨论过的，一个交易系统本质上是由多个交易子系统组成的投资组合，而每个品种对应一个子系统。你需要决定在这个组合内部如何分配品种权重。由于你使用的是静态交易规则，因此确定这些权重就成了一个关键决策。我在第四章中介绍了两种投资组合配置方法。Bootstrapping 会是很好的做法，但在这个例子里，我仍然采用更简单的手工构造法，并使用经验法则相关性，而不是回测得到的相关性。\end{chineseblock}

\bipairgap

\begin{englishblock} Remember from chapter eleven that as asset allocating investors with a static, rather than dynamic, trading strategy you can use the unadjusted correlation of instrument returns as a proxy for trading subsystem returns. Table 40 shows some indicative correlations, based mainly on the tables of indicative asset return correlations in appendix C, without any adjustment. As I discussed at the beginning of the chapter, you also have some constraints on your allocations. I'll explain in the optimisation where these affect your weights. \end{englishblock}

\begin{chineseblock}请记得，在第十一章中我提到过：由于资产配置型投资者使用的是静态而不是动态策略，因此你可以直接用未调整的品种收益相关性，作为交易子系统收益相关性的代理。表 40 展示了一些经验相关性，主要来自附录 C 中那些资产收益相关性的经验表格，并且没有做任何调整。正如本章开头所说，你的配置还受到一些额外约束。我会在下面的优化过程中解释，这些约束是如何影响权重的。\end{chineseblock}

\begin{bilingualtableblock} \rawtablecaption{WHAT ARE THE CORRELATIONS OF THE ETF PORTFOLIO?\\ETF 投资组合的相关性是多少？}\begingroup \scriptsize \setlength{\tabcolsep}{3pt} \renewcommand{\arraystretch}{1.08}
\par\noindent\makebox[\textwidth][l]{\resizebox{\textwidth}{!}{%
\begin{tabular}{@{} lcccccccccc@{}}
\toprule
& \shortstack{UK\\Bo.} & \shortstack{US\\Bo.} & \shortstack{EU\\Bo.} & \shortstack{EM\\Bo.} & \shortstack{Inf.\\Bo.} & \shortstack{EU\\Eq.} & \shortstack{US\\Eq.} & \shortstack{UK\\Eq.} & \shortstack{JP\\Eq.} & \shortstack{EM\\Eq.} \\
& \shortstack{英\\债} & \shortstack{美\\债} & \shortstack{欧\\债} & \shortstack{新兴\\市场债} & \shortstack{通胀\\挂钩债} & \shortstack{欧\\股} & \shortstack{美\\股} & \shortstack{英\\股} & \shortstack{日\\股} & \shortstack{新兴\\市场股} \\
\midrule
UK bond & 1 & & & & & & & & & \\
英国债券 & 1 & & & & & & & & & \\
US bond & 0.75 & 1 & & & & & & & & \\
美国债券 & 0.75 & 1 & & & & & & & & \\
EU bond & 0.75 & 0.75 & 1 & & & & & & & \\
欧洲债券 & 0.75 & 0.75 & 1 & & & & & & & \\
EM bond & 0.35 & 0.35 & 0.35 & 1 & & & & & & \\
新兴市场债券 & 0.35 & 0.35 & 0.35 & 1 & & & & & & \\
Infl. bond & 0.3 & 0.3 & 0.3 & 0.25 & 1 & & & & & \\
通胀挂钩债券 & 0.3 & 0.3 & 0.3 & 0.25 & 1 & & & & & \\
Euro Eq. & 0.1 & 0.1 & 0.1 & 0.1 & 0.1 & 1 & & & & \\
欧洲股票 & 0.1 & 0.1 & 0.1 & 0.1 & 0.1 & 1 & & & & \\
US equity & 0.1 & 0.1 & 0.1 & 0.1 & 0.1 & 0.75 & 1 & & & \\
美国股票 & 0.1 & 0.1 & 0.1 & 0.1 & 0.1 & 0.75 & 1 & & & \\
UK equity & 0.1 & 0.1 & 0.1 & 0.1 & 0.1 & 0.75 & 0.75 & 1 & & \\
英国股票 & 0.1 & 0.1 & 0.1 & 0.1 & 0.1 & 0.75 & 0.75 & 1 & & \\
JP equity & 0.1 & 0.1 & 0.1 & 0.1 & 0.1 & 0.75 & 0.75 & 0.75 & 1 & \\
日本股票 & 0.1 & 0.1 & 0.1 & 0.1 & 0.1 & 0.75 & 0.75 & 0.75 & 1 & \\
EM equity & 0.1 & 0.1 & 0.1 & 0.1 & 0.1 & 0.5 & 0.5 & 0.5 & 0.5 & 1 \\
新兴市场股票 & 0.1 & 0.1 & 0.1 & 0.1 & 0.1 & 0.5 & 0.5 & 0.5 & 0.5 & 1 \\
\bottomrule
\end{tabular}%
}}\endgroup \end{bilingualtableblock}

\begin{englishblock} The table shows the correlation matrix of instrument returns for the asset allocating example. Figures from tables 50 to 54, and my own estimates. We don’t need to adjust instrument returns when calculating instrument weights, since we’re trading a static portfolio. \end{englishblock}

\begin{chineseblock}这张表展示的是资产配置型示例中各品种收益的相关矩阵。相关性数值来自表 50 到表 54，再加上一部分我自己的估计。由于我们交易的是静态投资组合，因此在计算品种权重时不需要对品种收益做任何调整。\end{chineseblock}

\bipairgap

\begin{englishblock} Here is how I handcrafted the portfolio weights. Row numbers given refer to table 8 (page 79). \end{englishblock}

\begin{chineseblock}下面是我如何手工构造这组投资组合权重。文中提到的行号对应的是表 8（第 79 页）。\end{chineseblock}

\begin{bilingualtableblock} \noindent{\small \setlength{\tabcolsep}{2pt} \renewcommand{\arraystretch}{1.15} \begin{tabularx}{\textwidth}{@{} L{0.24\textwidth} Y@{}}\textbf{First level grouping} \par Within region/sub class &Developed bonds: 33.3\% to each of US, UK, Europe. Row 3. \par Developed equities: 25\% to each of US, UK, Europe, Japan. Row 3. \par All other groups have a single asset, with 100\% allocation. Row 1. \\

\textbf{第一层分组} \par 区域/子类别内部 &发达市场债券：美国、英国、欧洲各 33.3\%。对应第 3 行。 \par 发达市场股票：美国、英国、欧洲、日本各 25\%。对应第 3 行。 \par 其余各组都只有单一资产，因此分配 100\%。对应第 1 行。 \\
\addlinespace[0.6em]
\textbf{Second level grouping} \par Within asset class &Bonds: 25\% to emerging markets (constrained) and 25\% to inflation linked (constrained); leaves 50\% left over for developed markets. \par Equities: 30\% to emerging markets (constrained), leaves 70\% to developed markets. \\
债券：新兴市场债券 25\%（受约束），通胀挂钩债券 25\%（受约束），剩下的 50\% 分配给发达市场债券。 \par 股票：新兴市场股票 30\%（受约束），其余 70\% 分配给发达市场股票。 \\
\addlinespace[0.6em]
\textbf{Top level grouping} \par Across asset classes &40\% in bonds (constrained), leaves 60\% in equities. \\
债券占 40\%（受约束），其余 60\% 分配给股票。 \\
\end{tabularx}
}
\end{bilingualtableblock}

\begin{englishblock} You can see the final weights in table 41. Using these weights, the correlations and the formula on page 297, I get an instrument diversification multiplier of 1.61. \end{englishblock}

\begin{chineseblock}你可以在表 41 中看到最终权重。使用这些权重、相关性以及第 297 页的公式，我算出的品种分散化乘数为 1.61。 \end{chineseblock}

\begin{bilingualtableblock} \rawtablecaption{WHAT INSTRUMENT WEIGHTS SHOULD YOU USE FOR THE ETF PORTFOLIO?\\ETF 投资组合应使用怎样的品种权重？}\noindent{\small \setlength{\tabcolsep}{4pt} \renewcommand{\arraystretch}{1.12} \begin{tabularx}{\textwidth}{@{} p{0.22\textwidth} p{0.16\textwidth} p{0.16\textwidth} p{0.16\textwidth} Y@{}}
\toprule
\textbf{Instrument} & \textbf{Within region/sub class} & \textbf{Within asset class} & \textbf{Across asset classes} & \textbf{Final weight} \\
\textbf{品种} & \textbf{区域/子类别内权重} & \textbf{资产类别内权重} & \textbf{跨资产类别权重} & \textbf{最终权重} \\
\midrule
US bonds & 33.3\% & 50\% & 40\% & 6.67\% \\
美国债券 & 33.3\% & 50\% & 40\% & 6.67\% \\
Euro bonds & 33.3\% & 50\% & 40\% & 6.67\% \\
欧元区债券 & 33.3\% & 50\% & 40\% & 6.67\% \\
UK bonds & 33.3\% & 50\% & 40\% & 6.67\% \\
英国债券 & 33.3\% & 50\% & 40\% & 6.67\% \\
EM bonds & 100\% & 25\% & 40\% & 10\% \\
新兴市场债券 & 100\% & 25\% & 40\% & 10\% \\
Inflation bonds & 100\% & 25\% & 40\% & 10\% \\
通胀挂钩债券 & 100\% & 25\% & 40\% & 10\% \\
Euro Stoxx 50 & 25\% & 70\% & 60\% & 10.5\% \\
Euro Stoxx 50 & 25\% & 70\% & 60\% & 10.5\% \\
S\&P 500 & 25\% & 70\% & 60\% & 10.5\% \\
S\&P 500 & 25\% & 70\% & 60\% & 10.5\% \\
UK equities & 25\% & 70\% & 60\% & 10.5\% \\
英国股票 & 25\% & 70\% & 60\% & 10.5\% \\
Japan equities & 25\% & 70\% & 60\% & 10.5\% \\
日本股票 & 25\% & 70\% & 60\% & 10.5\% \\
EM equities & 100\% & 30\% & 60\% & 18\% \\
新兴市场股票 & 100\% & 30\% & 60\% & 18\% \\
\bottomrule
\end{tabularx}
}
\end{bilingualtableblock}

\begin{englishblock} The table shows the handcrafted instrument weights for the asset allocating investor. Final weights are the product of group weights at each stage. Weights in bold are constrained by the institutional mandate. Borders show groups. \end{englishblock}

\begin{chineseblock}这张表展示的是资产配置型投资者使用的手工构造品种权重。最终权重等于每个阶段分组权重的乘积。加粗的权重受到机构投资授权约束。边框则表示不同分组。\end{chineseblock}

\bilingualsubsection{Portfolios: Using predictions of performance}{投资组合：如何利用对表现的预测}

\begin{englishblock} Sometimes you might want to incorporate views about asset returns into investment portfolios. In an institutional setting you'll often have internal or external opinions on asset returns like "The Euro Stoxx will hit 3,000 in 12 months time." To use these figures, they first need to be translated into annualised Sharpe ratios (SR) using your current estimate of price volatility. A 3,000 Euro Stoxx target is about an 8\% rise from the price as I write this. With annualised price volatility of 16\% (derived from the daily value of 1\%), this equates to an annualised Sharpe ratio of 8\% ÷ 16\% = 0.50. You could do this within the framework by using discretionary forecasts, which would make this example look more like the semi-automatic trader of the previous chapter. However, you'd end up trading a lot more, which doesn't fit well with your objectives as a long run asset allocating investor. Instead I would recommend changing the handcrafted instrument weights within the portfolio using the SR weight adjustments from table 12 in chapter four (page 86), according to any opinions you have about performance. \end{englishblock}

\begin{chineseblock}有时候，你可能想把自己对资产收益的看法纳入投资组合。机构投资环境中，经常会出现内部或外部对资产收益的判断，例如：“12 个月后 Euro Stoxx 会到 3,000 点。”要把这类判断真正用起来，它们首先必须依据当前对价格波动率的估计，被转换成年化夏普比率（SR）。以我写作时的价格为例，Euro Stoxx 到 3,000 点大约意味着上涨 8\%。如果其年化价格波动率是 16\%（由日度 1\% 推导而来），那么这就对应于一个年化夏普比率：8\% ÷ 16\% = 0.50。你当然也可以在框架内部通过主观预测值来处理这件事，但那样一来，这个例子就会更像上一章里的半自动交易者。可这样做会导致你交易得多得多，而这并不符合长期资产配置投资者的目标。因此，我更建议你根据自己对各资产表现的判断，使用第四章表 12（第 86 页）中的 SR 权重调整方式，直接去调整投资组合里的手工构造品种权重。\end{chineseblock}

\bipairgap

\begin{englishblock} You should use column B ‘Without certainty' in the table, which reflects the difficulty in guessing how assets will perform. As an example, if the average SR across your instruments was 0.20, and you expected the Euro Stoxx SR to be 0.5, then with an expected 0.3 SR outperformance you'd increase your Euro Stoxx instrument weight by 17\%. Don't forget to renormalise the weights to add up to 100\% after adjustment. Please note that you don't need to change your instrument diversification multiplier when adjusting instrument weights, as this is designed to get your long run risk target correct. Finally, these adjustments should ideally be small and infrequent, since they are a source of additional trading and so will increase costs. \end{englishblock}

\begin{chineseblock}你应当使用该表中的 B 列“Without certainty”，因为它能够反映出：对资产表现做判断本身有多困难。举个例子，如果你所有品种的平均 SR 为 0.20，而你预计 Euro Stoxx 的 SR 会达到 0.5，那么它相对于平均值就有 0.3 的 SR 超额表现，于是你应当把 Euro Stoxx 的品种权重提高 17\%。别忘了，调整完之后还需要把所有权重重新归一化，使其重新加总为 100\%。同时请注意，在调整品种权重时，你并不需要修改品种分散化乘数，因为它的设计目的，是为了保证你的长期风险目标保持正确。最后，这类调整理想情况下应当既小又不频繁，因为它们本身会引发额外交易，从而增加成本。\end{chineseblock}

\bilingualsection{Weekly process}{每周流程}

\begin{englishblock} I've suggested a weekly rebalancing process here which given the low volatility target, lack of leverage and slow turnover should be fine. In practice rebalancing could be done more frequently for large funds or in times of market stress, others may choose to rebalance quarterly after meetings of investment committees, and amateur investors might be happy with annual rebalancing at the end of each tax year. Detailed calculations are shown in the trading diary section which follows. \end{englishblock}

\begin{chineseblock}在这个例子里，我建议使用每周再平衡流程。考虑到这里的波动率目标较低、不使用杠杆，而且换手很慢，这种频率完全足够。在实际中，对于大型基金或市场压力较大时期，再平衡也可以做得更频繁；另一些机构则可能选择在投资委员会开会之后，按季度再平衡；而普通投资者可能会满足于每个税务年度末进行一次年度再平衡。下面交易日记部分会展示详细计算。\end{chineseblock}

\begin{bilingualtableblock} \noindent{\small \setlength{\tabcolsep}{2pt} \renewcommand{\arraystretch}{1.15} \begin{tabularx}{\textwidth}{@{} p{0.29\textwidth} Y@{}}\textbf{Get account value} & Get today's account value and work out current capital. \\

\textbf{获取账户净值} & 获取当天账户净值，并计算当前资本。 \\
\textbf{Cash volatility target} & Your annualised cash volatility target equals the percentage volatility target multiplied by current capital. In this example we're using a 7.5\% percentage volatility target. Divide the annualised target by 16 for daily cash volatility target. \\
\textbf{现金波动率目标} & 年化现金波动率目标等于百分比波动率目标乘以当前资本。本例中我们使用 7.5\% 的百分比波动率目标。将年化目标除以 16，就得到日现金波动率目标。 \\
\textbf{Get latest prices} & Get prices for all instruments. \\
\textbf{获取最新价格} & 获取所有品种的最新价格。 \\
\textbf{Get latest FX rates} & Update FX rates. For this example you only need GBP/EUR and USD/EUR. \\
\textbf{获取最新汇率} & 更新汇率。本例中你只需要 GBP/EUR 和 USD/EUR。 \\
\textbf{Calculate instrument value volatility} & Using the recommended 20 week moving average look-back, work out the price volatility of each instrument. With the FX rates and block value convert this to instrument value volatility. \\
\textbf{计算品种价值波动率} & 使用建议的 20 周移动平均回看窗口，算出每个品种的价格波动率。再结合汇率和单位价值，把它转换为品种价值波动率。 \\
\textbf{Instrument subsystem position} & Because you have a constant forecast of +10 this is equal to the volatility scalar, that is daily cash volatility target divided by instrument value volatility. \\
\textbf{品种子系统仓位} & 由于你的预测值恒定为 +10，所以这一步的结果就等于波动率缩放因子，也就是日现金波动率目标除以品种价值波动率。 \\
\textbf{Optional: Get estimate of performance} & Translate predictions into annualised Sharpe ratios using your predictions for price volatility. \\
\textbf{可选：获取表现预估} & 把你对价格表现的判断转换成年化夏普比率。 \\
\textbf{Optional: Adjust instrument weights} & Using Sharpe ratio predictions table 12 (page 86), column B ‘Without certainty'. \\
\textbf{可选：调整品种权重} & 使用第四章表 12（第 86 页）B 列 “Without certainty” 来进行夏普比率权重调整。 \\
\textbf{Portfolio instrument position} & Subsystem position multiplied by instrument weight, from table 41 for the example, with Sharpe ratio adjustments if required, and by instrument diversification multiplier, 1.61 in the example. \\
\textbf{投资组合品种仓位} & 子系统仓位乘以品种权重（本例使用表 41 的权重，如有需要再做夏普比率调整），并再乘以品种分散化乘数（本例中为 1.61）。 \\
\textbf{Rounded target position} & Portfolio instrument position rounded to nearest instrument block position, in the example 100 shares. \\
\textbf{四舍五入后的目标仓位} & 将投资组合品种仓位四舍五入到最接近的品种单位仓位，本例中一个单位是 100 股。 \\
\textbf{Issue trades} & Compare current position to final desired position. If out by more than 10\% then trade, using position inertia. \\
\textbf{发出交易指令} & 把当前仓位与最终目标仓位进行比较。如果偏差超过 10\%，就依据仓位惯性进行交易。 \\
\end{tabularx}
}
\end{bilingualtableblock}

\bilingualsection{Trading diary}{投资日记}

\begin{englishblock} This is a hypothetical and contrived example which I've constructed to show the main features of a system that could have been trading during the great crash of 2008.\footnote{Although I do run a similar portfolio to this I have only done so since 2011. Most of the time very little happens in a portfolio like this, so I've decided to set this example in a more interesting historical period. \par 虽然我自己确实运行着一个相似的投资组合，但那是从 2011 年才开始的。像这种组合在大多数时间里其实都很平淡，所以我决定把这个例子放在一个更有戏剧性的历史时期里。} Since many of the ETFs being traded did not exist at the time I am using today's prices, although the subsequent price changes and volatility levels are in line with what you would have seen in the past. As they are not a key part of the example I'm also using arbitrary fixed exchange rates. Spreadsheets detailing all the calculations are available on the website for this book, \url{www.systematictrading.org}. \end{englishblock}

\begin{chineseblock}这是一个假设性的、经过人为构造的例子，我用它来展示：如果你在 2008 年大崩盘期间运行这样一套系统，它会呈现出怎样的主要特征。由于这里交易的许多 ETF 在当时并不存在，因此我使用的是今天的价格，只不过后续的价格变化与波动率水平，被设定得与当时你真实会看到的情况相匹配。由于汇率并不是这个例子的关键部分，我也在这里使用了一组人为固定的汇率。所有详细计算的电子表格，都可以在本书网站\url{www.systematictrading.org}上找到。\end{chineseblock}

\bilingualsubsection{1 July 2008}{2008 年 7 月 1 日}

\begin{englishblock} You begin with trading capital of €10,000,000, which with your 7.5\% volatility target gives a daily volatility target of one sixteenth of €750,000, or €46,875. You haven't yet got any discretionary forecasts for asset prices, so you can keep the default unadjusted instrument weights. \end{englishblock}

\begin{chineseblock}你一开始拥有 €10,000,000 的交易资本，而在 7.5\% 波动率目标下，对应的日波动率目标就是 €750,000 的十六分之一，也就是 €46,875。由于你此时还没有引入任何关于资产价格的主观预测，因此可以直接使用默认的、未经调整的品种权重。\end{chineseblock}

\bilingualsubsection{Bond Calculations Part One}{债券部分计算：第一部分}

\begin{chinesetableblock} \noindent{\scriptsize \setlength{\tabcolsep}{2pt} \renewcommand{\arraystretch}{1.12} \begin{tabularx}{\textwidth}{@{} p{0.30\textwidth}*{5}{R}@{}}
\toprule
& \$ US bonds / 美元美国债券 & € Euro bonds / 欧元欧洲债券 & £ UK bonds / 英镑英国债券 & \$ EM bonds / 美元新兴市场债券 & \$ Inflation bonds / 美元通胀挂钩债券 \\
\midrule
Daily volatility target / 日波动率目标 (A) & €46,875 & €46,875 & €46,875 & €46,875 & €46,875 \\
Price / 价格 (B) & \$5 & €150 & £22 & \$75 & \$145 \\
Number of shares per block / 每单位股数 (C) & 100 & 100 & 100 & 100 & 100 \\
FX / 汇率 (D) & 0.9 & 1 & 1.2 & 0.9 & 0.9 \\
Block value / 单位价值 (E) & \$5 & €150 & £22 & \$75 & \$145 \\
\emph{C × B ÷ 100} & & & & & \\
Price volatility / 价格波动率, \% (G) & 0.90 & 0.60 & 0.55 & 0.70 & 0.43 \\
Instrument currency volatility / 品种货币波动率 (H) & \$4.5 & €90 & £12.1 & \$52.5 & \$62.35 \\
\emph{G × E} & & & & & \\
Instrument value volatility / 品种价值波动率 (I) & €4.05 & €90.00 & €14.52 & €47.25 & €56.12 \\
\emph{H × D} & & & & & \\
Volatility scalar / 波动率缩放因子 (J) & 11,574 & 521 & 3228 & 992 & 835 \\
\emph{A ÷ I} & & & & & \\
\bottomrule
\end{tabularx}
}
\end{chinesetableblock}

\smallskip \noindent\textit{I've kept the row identifiers the same in all three example chapters so that comparisons can easily be made. Row F has been omitted as you don't require volatility in points per day here.}

\smallskip \noindent\textit{我在三个示例章节中都保留了相同的行标识，便于直接比较。由于这里不需要“以点数表示的日波动率”，所以省略了 F 行。}

\medskip

\bilingualsubsection{Bond Calculations Part Two}{债券部分计算：第二部分}

\begin{chinesetableblock} \noindent{\scriptsize \setlength{\tabcolsep}{2pt} \renewcommand{\arraystretch}{1.12} \begin{tabularx}{\textwidth}{@{} p{0.30\textwidth}*{5}{R}@{}}
\toprule
& \$ US bonds / 美元美国债券 & € Euro bonds / 欧元欧洲债券 & £ UK bonds / 英镑英国债券 & \$ EM bonds / 美元新兴市场债券 & \$ Inflation bonds / 美元通胀挂钩债券 \\
\midrule
Forecast / 预测值 (K) & +10 & +10 & +10 & +10 & +10 \\
Subsystem position, blocks / 子系统仓位，单位 (L) & 11574 & 521 & 3228 & 992 & 835 \\
\emph{K × J ÷ 10} & & & & & \\
Instrument weight / 品种权重 (M) & 6.67\% & 6.67\% & 6.67\% & 10.00\% & 10.00\% \\
Instrument diversification multiplier / 品种分散化乘数 (N) & 1.61 & 1.61 & 1.61 & 1.61 & 1.61 \\
Portfolio instrument position, blocks / 投资组合品种仓位，单位 (O) & 1242.3 & 55.9 & 346.5 & 159.7 & 134.5 \\
\emph{L × M × N} & & & & & \\
Rounded target position, blocks / 四舍五入后的目标仓位，单位 (P = round O) & 1242 & 56 & 347 & 160 & 135 \\
\bottomrule
\end{tabularx}
}
\end{chinesetableblock}

\begin{englishblock} So with 100 share blocks you buy 1242 × 100 = 124,200 shares of the US bond ETF and so on. \end{englishblock}

\begin{chineseblock}因此，在每个交易单位为 100 股的设定下，你会买入 1242 × 100 = 124,200 股美国债券 ETF，其他品种以此类推。\end{chineseblock}

\bilingualsubsection{Equity Calculations Part One}{股票部分计算：第一部分}

\begin{chinesetableblock} \noindent{\scriptsize \setlength{\tabcolsep}{2pt} \renewcommand{\arraystretch}{1.12} \begin{tabularx}{\textwidth}{@{} p{0.30\textwidth}*{5}{R}@{}}
\toprule
& € Euro equities / 欧元欧洲股票 & \$ US equities / 美元美国股票 & £ UK equities / 英镑英国股票 & € Japan equities / 欧元日本股票 & \$ EM equities / 美元新兴市场股票 \\
\midrule
Daily volatility target / 日波动率目标 (A) & €46,875 & €46,875 & €46,875 & €46,875 & €46,875 \\
Price / 价格 (B) & €37 & \$40 & £7 & €15 & \$54 \\
Number of shares per block / 每单位股数 (C) & 100 & 100 & 100 & 100 & 100 \\
FX / 汇率 (D) & 1 & 0.9 & 1.2 & 1 & 0.9 \\
Block value / 单位价值 (E) & €37 & \$40 & £7 & €15 & \$54 \\
\emph{C × B ÷ 100} & & & & & \\
Price volatility / 价格波动率, \% (G) & 1.10 & 1.10 & 1.20 & 1.00 & 1.40 \\
Instrument currency volatility / 品种货币波动率 (H) & €40.7 & \$44 & £8.4 & €15 & \$75.6 \\
\emph{G × E} & & & & & \\
Instrument value volatility / 品种价值波动率 (I) & €40.70 & €39.60 & €10.08 & €15.00 & €68.04 \\
\emph{H × D} & & & & & \\
Volatility scalar / 波动率缩放因子 (J) & 1152 & 1184 & 4650 & 3125 & 689 \\
\emph{A ÷ I} & & & & & \\
\bottomrule
\end{tabularx}
}
\end{chinesetableblock}

\bilingualsubsection{Equity Calculations Part Two}{股票部分计算：第二部分}

\begin{chinesetableblock} \noindent{\scriptsize \setlength{\tabcolsep}{2pt} \renewcommand{\arraystretch}{1.12} \begin{tabularx}{\textwidth}{@{} p{0.30\textwidth}*{5}{R}@{}}
\toprule
& € Euro equities / 欧元欧洲股票 & \$ US equities / 美元美国股票 & £ UK equities / 英镑英国股票 & € Japan equities / 欧元日本股票 & \$ EM equities / 美元新兴市场股票 \\
\midrule
Forecast / 预测值 (K) & +10 & +10 & +10 & +10 & +10 \\
Subsystem position, blocks / 子系统仓位，单位 (L) & 1152 & 1184 & 4650 & 3125 & 689 \\
\emph{K × J ÷ 10} & & & & & \\
Instrument weight / 品种权重 (M) & 10.50\% & 10.50\% & 10.50\% & 10.50\% & 18.00\% \\
Instrument diversification multiplier / 品种分散化乘数 (N) & 1.61 & 1.61 & 1.61 & 1.61 & 1.61 \\
Portfolio instrument position, blocks / 投资组合品种仓位，单位 (O) & 194.7 & 200.1 & 786.1 & 528.3 & 199.7 \\
\emph{L × M × N} & & & & & \\
Rounded target position, and trade / 四舍五入后的目标仓位与交易 (P = round O) & 195 & 200 & 786 & 528 & 200 \\
\bottomrule
\end{tabularx}
}
\end{chinesetableblock}

\begin{englishblock} An interesting statistic is that 57\% of the cash value of this portfolio is in bonds, despite them accounting for only 40\% of the instrument weight. You need more in bonds because they are lower risk; this reflects the risk parity nature of the portfolio. \end{englishblock}

\begin{chineseblock}一个有意思的事实是：这个投资组合中有 57\% 的现金价值实际上配置在债券上，尽管债券在风险调整后的品种权重中只占 40\%。之所以债券的现金金额更多，是因为它们风险更低；这恰恰体现了这个组合的风险平价性质。\end{chineseblock}

\bilingualsubsection{1 October 2008}{2008 年 10 月 1 日}

\begin{englishblock} The world is getting riskier but despite the Lehman's bankruptcy a couple of weeks ago the portfolio hasn't yet seen significant losses so your current capital is almost unchanged at €9,705,000. Your economic strategists feel the fear is overblown and are (a) insanely bullish on equities with an expected Sharpe ratio (SR) of 0.6 and (b) expecting all bonds to make exactly nothing; an SR of zero. \end{englishblock}

\begin{chineseblock}世界正在变得越来越危险，但尽管几周前雷曼已经破产，这个投资组合目前还没有遭受特别明显的损失，因此你的当前资本几乎没怎么变化，仍然大约是 €9,705,000。你的经济策略团队认为市场恐慌被夸大了，因此（a）他们对股票极度看多，预期夏普比率（SR）达到 0.6； （b）他们则认为所有债券都不会带来任何收益，也就是 SR 为零。\end{chineseblock}

\bipairgap

\begin{englishblock} Using these predictions you can use table 12 (page 86), column B ‘Without certainty', to calculate which Sharpe ratio adjustments to make. Since there is no performance differential within handcrafted groups we can do the adjustments in one go, treating the entire portfolio as a single group. The simple average of portfolio performance is the average of 0, for bonds, and 0.6, for equities, or an SR of 0.3. \end{englishblock}

\begin{chineseblock}利用这些判断，你可以使用表 12（第 86 页）中的 B 列“Without certainty”，来计算应该进行哪些夏普比率调整。由于在手工构造分组内部，我们并不区分各成员之间的表现差异，因此可以把整个投资组合视作一个整体，一次性完成调整。组合表现的简单平均值，就是债券的 0 和股票的 0.6 取平均，也就是 SR = 0.3。 \end{chineseblock}

\bilingualsubsection{Bond Instrument Weight Adjustments}{债券品种权重调整}

\begin{chinesetableblock} \noindent{\scriptsize \setlength{\tabcolsep}{2pt} \renewcommand{\arraystretch}{1.12} \begin{tabularx}{\textwidth}{@{} p{0.30\textwidth}*{5}{R}@{}}
\toprule
& US bonds / 美国债券 & Euro bonds / 欧洲债券 & UK bonds / 英国债券 & EM bonds / 新兴市场债券 & Inflation bonds / 通胀挂钩债券 \\
\midrule
Instrument weight / 品种权重 & 6.67\% & 6.67\% & 6.67\% & 10.00\% & 10.00\% \\
Expected Sharpe ratio / 预期夏普比率 & 0.0 & 0.0 & 0.0 & 0.0 & 0.0 \\
Average Sharpe ratio / 平均夏普比率 & 0.3 & 0.3 & 0.3 & 0.3 & 0.3 \\
Outperformance / 相对超额表现 & -0.3 & -0.3 & -0.3 & -0.3 & -0.3 \\
Sharpe ratio adjustment / 夏普比率调整* & 0.83 & 0.83 & 0.83 & 0.83 & 0.83 \\
Renormalised instrument weight / 重新归一化后的品种权重** & 5.35\% & 5.35\% & 5.35\% & 8.03\% & 8.03\% \\
\bottomrule
\end{tabularx}
}
\end{chinesetableblock}

\smallskip \noindent\textit{* From table 12 (page 86) column B. ** After adjustment, weights add up to 103.4\%. So we need to divide all adjusted weights by 1.034 to ensure they sum to 100\%.}

\smallskip \noindent\textit{* 来自表 12（第 86 页）B 列。** 调整后，所有权重加总为 103.4\%，因此必须再把所有调整后的权重都除以 1.034，才能让它们重新加总为 100\%。}

\medskip

\bilingualsubsection{Equity Instrument Weight Adjustments}{股票品种权重调整}

\begin{chinesetableblock} \noindent{\scriptsize \setlength{\tabcolsep}{2pt} \renewcommand{\arraystretch}{1.12} \begin{tabularx}{\textwidth}{@{} p{0.30\textwidth}*{5}{R}@{}}
\toprule
& Euro equities / 欧洲股票 & US equities / 美国股票 & UK equities / 英国股票 & Japan equities / 日本股票 & EM equities / 新兴市场股票 \\
\midrule
Instrument weight / 品种权重 & 10.50\% & 10.50\% & 10.50\% & 10.50\% & 18.00\% \\
Expected Sharpe ratio / 预期夏普比率 & 0.6 & 0.6 & 0.6 & 0.6 & 0.6 \\
Average Sharpe ratio / 平均夏普比率 & 0.3 & 0.3 & 0.3 & 0.3 & 0.3 \\
Outperformance / 相对超额表现 & 0.3 & 0.3 & 0.3 & 0.3 & 0.3 \\
Sharpe ratio adjustment / 夏普比率调整 & 1.17 & 1.17 & 1.17 & 1.17 & 1.17 \\
Renormalised instrument weight / 重新归一化后的品种权重 & 11.88\% & 11.88\% & 11.88\% & 11.88\% & 20.37\% \\
\bottomrule
\end{tabularx}
}
\end{chinesetableblock}

\bilingualsubsection{Bond Calculations Part One}{债券部分计算：第一部分（更新后）}

\begin{chinesetableblock} \noindent{\scriptsize \setlength{\tabcolsep}{2pt} \renewcommand{\arraystretch}{1.12} \begin{tabularx}{\textwidth}{@{} p{0.30\textwidth}*{5}{R}@{}}
\toprule
& \$ US bonds / 美元美国债券 & € Euro bonds / 欧元欧洲债券 & £ UK bonds / 英镑英国债券 & \$ EM bonds / 美元新兴市场债券 & \$ Inflation bonds / 美元通胀挂钩债券 \\
\midrule
Daily volatility target / 日波动率目标 (A) & €45,492 & €45,492 & €45,492 & €45,492 & €45,492 \\
Price / 价格 (B) & \$5.10 & €160 & £22.15 & \$72 & \$140 \\
Number of shares per block / 每单位股数 (C) & 100 & 100 & 100 & 100 & 100 \\
FX / 汇率 (D) & 0.9 & 1 & 1.2 & 0.9 & 0.9 \\
Block value / 单位价值 (E) & \$5.10 & €160 & £22.15 & \$72 & \$140 \\
\emph{C × B ÷ 100} & & & & & \\
Price volatility / 价格波动率, \% (G) & 1.00 & 0.65 & 0.55 & 1.00 & 0.46 \\
Instrument currency volatility / 品种货币波动率 (H) & \$5.1 & €104 & £12.1825 & \$72 & \$64.4 \\
\emph{G × E} & & & & & \\
Instrument value volatility / 品种价值波动率 (I) & €4.59 & €104.00 & €14.62 & €64.80 & €57.96 \\
\emph{H × D} & & & & & \\
Volatility scalar / 波动率缩放因子 (J) & 9911 & 437 & 3112 & 702 & 785 \\
\emph{A ÷ I} & & & & & \\
\bottomrule
\end{tabularx}
}
\end{chinesetableblock}

\smallskip \noindent\textit{* I've kept the row identifiers the same in all three example chapters so that comparisons can easily be made. Row F has been omitted as you don't require volatility in points per day here.}

\smallskip \noindent\textit{* 我在三个示例章节中都保留了相同的行标识，便于直接比较。由于这里不需要“以点数表示的日波动率”，因此省略了 F 行。}

\medskip

\bilingualsubsection{Bond Calculations Part Two}{债券部分计算：第二部分（更新后）}

\begin{chinesetableblock} \noindent{\scriptsize \setlength{\tabcolsep}{2pt} \renewcommand{\arraystretch}{1.12} \begin{tabularx}{\textwidth}{@{} p{0.30\textwidth}*{5}{R}@{}}
\toprule
& \$ US bonds / 美元美国债券 & € Euro bonds / 欧元欧洲债券 & £ UK bonds / 英镑英国债券 & \$ EM bonds / 美元新兴市场债券 & \$ Inflation bonds / 美元通胀挂钩债券 \\
\midrule
Forecast / 预测值 (K) & +10 & +10 & +10 & +10 & +10 \\
Subsystem position, blocks / 子系统仓位，单位 (L) & 9911 & 437 & 3112 & 702 & 785 \\
\emph{K × J ÷ 10} & & & & & \\
Instrument weight / 品种权重 (M) & 5.35\% & 5.35\% & 5.35\% & 8.03\% & 8.03\% \\
Instrument diversification multiplier / 品种分散化乘数 (N) & 1.61 & 1.61 & 1.61 & 1.61 & 1.61 \\
Portfolio instrument position, blocks / 投资组合品种仓位，单位 (O) & 853.9 & 37.7 & 268.1 & 90.7 & 101.4 \\
\emph{L × M × N} & & & & & \\
Rounded target position, blocks / 四舍五入后的目标仓位，单位 (P = round O) & 854 & 38 & 268 & 91 & 101 \\
Current position, blocks / 当前仓位，单位 & 1242 & 56 & 347 & 160 & 134 \\
Trade, blocks / 交易，单位 & Sell 388 / 卖出 388 & Sell 18 / 卖出 18 & Sell 79 / 卖出 79 & Sell 69 / 卖出 69 & Sell 33 / 卖出 33 \\
\bottomrule
\end{tabularx}
}
\end{chinesetableblock}

\bilingualsubsection{Equity Calculations Part One}{股票部分计算：第一部分（更新后）}

\begin{chinesetableblock} \noindent{\scriptsize \setlength{\tabcolsep}{2pt} \renewcommand{\arraystretch}{1.12} \begin{tabularx}{\textwidth}{@{} p{0.30\textwidth}*{5}{R}@{}}
\toprule
& € Euro equities / 欧元欧洲股票 & \$ US equities / 美元美国股票 & £ UK equities / 英镑英国股票 & ¥ Japan equities / 日元日本股票 & \$ EM equities / 美元新兴市场股票 \\
\midrule
Daily volatility target / 日波动率目标 (A) & €45,492 & €45,492 & €45,492 & €45,492 & €45,492 \\
Price / 价格 (B) & €35 & \$37 & £6.5 & €14 & \$50 \\
Number of shares per block / 每单位股数 (C) & 100 & 100 & 100 & 100 & 100 \\
FX / 汇率 (D) & 1 & 0.9 & 1.2 & 1 & 0.9 \\
Block value / 单位价值 (E) & €35 & \$37 & £6.5 & €14 & \$50 \\
\emph{C × B ÷ 100} & & & & & \\
Price volatility / 价格波动率, \% (G) & 1.20 & 1.30 & 1.30 & 1.10 & 1.60 \\
Instrument currency volatility / 品种货币波动率 (H) & €42 & \$48.1 & £8.45 & €15.4 & \$80 \\
\emph{G × E} & & & & & \\
Instrument value volatility / 品种价值波动率 (I) & €42 & €43.29 & €12 & €15.4 & €72 \\
\emph{H × D} & & & & & \\
Volatility scalar / 波动率缩放因子 (J) & 1083 & 1051 & 4486 & 2954 & 632 \\
\emph{A ÷ I} & & & & & \\
\bottomrule
\end{tabularx}
}
\end{chinesetableblock}

\bilingualsubsection{Equity Calculations Part Two}{股票部分计算：第二部分（更新后）}

\begin{chinesetableblock} \noindent{\scriptsize \setlength{\tabcolsep}{2pt} \renewcommand{\arraystretch}{1.12} \begin{tabularx}{\textwidth}{@{} p{0.30\textwidth}*{5}{R}@{}}
\toprule
& € Euro equities / 欧元欧洲股票 & \$ US equities / 美元美国股票 & £ UK equities / 英镑英国股票 & ¥ Japan equities / 日元日本股票 & \$ EM equities / 美元新兴市场股票 \\
\midrule
Forecast / 预测值 (K) & +10 & +10 & +10 & +10 & +10 \\
Subsystem position, blocks / 子系统仓位，单位 (L) & 1083 & 1051 & 4486 & 2954 & 632 \\
\emph{K × J ÷ 10} & & & & & \\
Instrument weight / 品种权重 (M) & 11.88\% & 11.88\% & 11.88\% & 11.88\% & 20.37\% \\
Instrument diversification multiplier / 品种分散化乘数 (N) & 1.61 & 1.61 & 1.61 & 1.61 & 1.61 \\
Portfolio instrument position, blocks / 投资组合品种仓位，单位 (O) & 207.2 & 201.0 & 858.2 & 565.1 & 207.2 \\
\emph{L × M × N} & & & & & \\
Rounded target position, blocks / 四舍五入后的目标仓位，单位 (P = round O) & 207 & 201 & 858 & 565 & 207 \\
Current position, blocks / 当前仓位，单位 & 195 & 200 & 786 & 528 & 200 \\
Trade, blocks / 交易，单位 & No / 无 & No / 无 & No / 无 & No / 无 & No / 无 \\
\bottomrule
\end{tabularx}
}
\end{chinesetableblock}

\begin{englishblock} You wouldn't trade any equity markets as all positions are within 10\% of their desired value. This is because the increase in equity price volatility, which would otherwise reduce positions, has mostly been compensated for by the increase in instrument weights. \end{englishblock}

\begin{chineseblock}在这一时点，你不会对任何股票市场进行交易，因为所有仓位都仍然处在目标值的 10\% 范围以内。这是因为，股票价格波动率的上升本来会压低仓位，但这一影响在很大程度上又被品种权重的上升所抵消了。\end{chineseblock}

\bilingualsubsection{10 October 2008}{2008 年 10 月 10 日}

\begin{englishblock} The equity markets fell off a cliff last week so you've come in on Saturday to assess the damage. Fortunately your portfolio is diversified enough that your losses are limited, even with the stupid equity overweight. You increased your equity allocation only modestly, from 60\% to 68\%, which reflects the uncertainty involved in forecasting returns. This means you just tilted the portfolio towards equities, rather than completely reallocating. But there has been some pain -- your current capital is now down to €9,126,500. \end{englishblock}

\begin{chineseblock}上周股票市场像坠崖一样暴跌，因此你周六也不得不来办公室评估损失。幸运的是，你的投资组合足够分散，因此即便你做了那个有点愚蠢的股票超配，总体损失依然受到控制。你只是把股票配置从 60\% 适度提高到了 68\%，这本身就反映出：对于收益预测，你仍然保留着相当程度的不确定性。也就是说，你只是把组合轻微地向股票倾斜，而不是进行彻底重配。不过，痛苦还是来了--- 你的当前资本已经下跌到 €9,126,500。 \end{chineseblock}

\bipairgap

\begin{englishblock} Even with your long look-back, price volatility has exploded, especially in equities. The economic strategists have belatedly decided that bonds are now the way to go and they have shifted to underweight in all equities. Your previous Sharpe ratio adjustments are now inverted. \end{englishblock}

\begin{chineseblock}即便你已经采用了很长的回看窗口，价格波动率依然出现了爆炸式上升，尤其是在股票市场。此时，经济策略团队也终于后知后觉地认为，债券才是当前更好的方向，于是他们转而对所有股票给出低配判断。你之前所做的夏普比率调整，现在也全部反了过来。\end{chineseblock}

\bilingualsubsection{Bond Instrument Weight Adjustments}{债券品种权重调整（反转后）}

\begin{chinesetableblock} \noindent{\scriptsize \setlength{\tabcolsep}{2pt} \renewcommand{\arraystretch}{1.12} \begin{tabularx}{\textwidth}{@{} p{0.30\textwidth}*{5}{R}@{}}
\toprule
& US bonds / 美国债券 & Euro bonds / 欧洲债券 & UK bonds / 英国债券 & EM bonds / 新兴市场债券 & Inflation bonds / 通胀挂钩债券 \\
\midrule
Instrument weight / 品种权重 & 6.67\% & 6.67\% & 6.67\% & 10.00\% & 10.00\% \\
Expected Sharpe ratio / 预期夏普比率 & 0.6 & 0.6 & 0.6 & 0.6 & 0.6 \\
Average Sharpe ratio / 平均夏普比率 & 0.3 & 0.3 & 0.3 & 0.3 & 0.3 \\
Outperformance / 相对超额表现 & 0.3 & 0.3 & 0.3 & 0.3 & 0.3 \\
Sharpe ratio adjustment / 夏普比率调整* & 1.17 & 1.17 & 1.17 & 1.17 & 1.17 \\
Renormalised instrument weight / 重新归一化后的品种权重** & 8.07\% & 8.07\% & 8.07\% & 12.11\% & 12.11\% \\
\bottomrule
\end{tabularx}
}
\end{chinesetableblock}

\smallskip \noindent\textit{* From table 12 (page 86) column B. ** After adjustment weights add up to 96.6\%. So you need to divide all adjusted weights by 0.966 to ensure they sum to 100\%.}

\smallskip \noindent\textit{* 来自表 12（第 86 页）B 列。** 调整后权重加总为 96.6\%，因此需要把所有调整后的权重都除以 0.966，以保证它们重新加总为 100\%。}

\medskip

\bilingualsubsection{Equity Instrument Weight Adjustments}{股票品种权重调整（反转后）}

\begin{chinesetableblock} \noindent{\scriptsize \setlength{\tabcolsep}{2pt} \renewcommand{\arraystretch}{1.12} \begin{tabularx}{\textwidth}{@{} p{0.30\textwidth}*{5}{R}@{}}
\toprule
& Euro equities / 欧洲股票 & US equities / 美国股票 & UK equities / 英国股票 & Japan equities / 日本股票 & EM equities / 新兴市场股票 \\
\midrule
Instrument weight / 品种权重 & 10.5\% & 10.5\% & 10.5\% & 10.5\% & 18\% \\
Expected Sharpe ratio / 预期夏普比率 & 0.0 & 0.0 & 0.0 & 0.0 & 0.0 \\
Average Sharpe ratio / 平均夏普比率 & 0.3 & 0.3 & 0.3 & 0.3 & 0.3 \\
Outperformance / 相对超额表现 & -0.3 & -0.3 & -0.3 & -0.3 & -0.3 \\
Sharpe ratio adjustment / 夏普比率调整 & 0.83 & 0.83 & 0.83 & 0.83 & 0.83 \\
Renormalised instrument weight / 重新归一化后的品种权重 & 9.02\% & 9.02\% & 9.02\% & 9.02\% & 15.47\% \\
\bottomrule
\end{tabularx}
}
\end{chinesetableblock}

\bilingualsubsection{Bond Calculations Part One}{债券部分计算：第一部分（再次更新）}

\begin{chinesetableblock} \noindent{\scriptsize \setlength{\tabcolsep}{2pt} \renewcommand{\arraystretch}{1.12} \begin{tabularx}{\textwidth}{@{} p{0.30\textwidth}*{5}{R}@{}}
\toprule
& \$ US bonds / 美元美国债券 & € Euro bonds / 欧元欧洲债券 & £ UK bonds / 英镑英国债券 & \$ EM bonds / 美元新兴市场债券 & \$ Inflation bonds / 美元通胀挂钩债券 \\
\midrule
Daily volatility target / 日波动率目标 (A) & €42,780.47 & €42,780.47 & €42,780.47 & €42,780.47 & €42,780.47 \\
Price / 价格 (B) & \$5.8 & €170 & £23.8 & \$70 & \$135 \\
Number of shares per block / 每单位股数 (C) & 100 & 100 & 100 & 100 & 100 \\
FX / 汇率 (D) & 0.9 & 1 & 1.2 & 0.9 & 0.9 \\
Block value / 单位价值 (E) & \$5.8 & €170 & £23.8 & \$70 & \$135 \\
\emph{C × B ÷ 100} & & & & & \\
Price volatility / 价格波动率, \% (G) & 1\% & 0.65\% & 0.6\% & 1.1\% & 0.6\% \\
Instrument currency volatility / 品种货币波动率 (H) & \$5.8 & €110.5 & £14.28 & \$77 & \$81 \\
\emph{G × E} & & & & & \\
Instrument value volatility / 品种价值波动率 (I) & €5.22 & €110.5 & €17.14 & €69.3 & €72.9 \\
\emph{H × D} & & & & & \\
Volatility scalar / 波动率缩放因子 (J) & 8195 & 387 & 2497 & 617 & 587 \\
\emph{A ÷ I} & & & & & \\
\bottomrule
\end{tabularx}
}
\end{chinesetableblock}

\smallskip \noindent\textit{* I've kept the row identifiers the same in all three example chapters so that comparisons can easily be made. Row F has been omitted as you don't require volatility in points per day here.}

\smallskip \noindent\textit{* 我在三个示例章节中保留了相同的行标识，便于直接比较。由于这里不需要“以点数表示的日波动率”，所以省略了 F 行。}

\medskip

\bilingualsubsection{Bond Calculations Part Two}{债券部分计算：第二部分（再次更新）}

\begin{chinesetableblock} \noindent{\scriptsize \setlength{\tabcolsep}{2pt} \renewcommand{\arraystretch}{1.12} \begin{tabularx}{\textwidth}{@{} p{0.30\textwidth}*{5}{R}@{}}
\toprule
& \$ US bonds / 美元美国债券 & € Euro bonds / 欧元欧洲债券 & £ UK bonds / 英镑英国债券 & \$ EM bonds / 美元新兴市场债券 & \$ Inflation bonds / 美元通胀挂钩债券 \\
\midrule
Forecast / 预测值 (K) & +10 & +10 & +10 & +10 & +10 \\
Subsystem position, blocks / 子系统仓位，单位 (L) & 8195 & 387 & 2497 & 617 & 587 \\
\emph{K × J ÷ 10} & & & & & \\
Instrument weight / 品种权重 (M) & 8.07\% & 8.07\% & 8.07\% & 12.11\% & 12.11\% \\
Instrument diversification multiplier / 品种分散化乘数 (N) & 1.61 & 1.61 & 1.61 & 1.61 & 1.61 \\
Portfolio instrument position, blocks / 投资组合品种仓位，单位 (O) & 1065.4 & 50.3 & 324.5 & 120.4 & 114.4 \\
\emph{L × M × N} & & & & & \\
Rounded target position, blocks / 四舍五入后的目标仓位，单位 (P = round O) & 1065 & 50 & 325 & 120 & 114 \\
Current position, blocks / 当前仓位，单位 & 854 & 38 & 268 & 91 & 101 \\
Trade, blocks / 交易，单位 & Buy 211 / 买入 211 & Buy 12 / 买入 12 & Buy 57 / 买入 57 & Buy 29 / 买入 29 & Buy 13 / 买入 13 \\
\bottomrule
\end{tabularx}
}
\end{chinesetableblock}

\bilingualsubsection{Equity Calculations Part One}{股票部分计算：第一部分（再次更新）}

\begin{chinesetableblock} \noindent{\scriptsize \setlength{\tabcolsep}{2pt} \renewcommand{\arraystretch}{1.12} \begin{tabularx}{\textwidth}{@{} p{0.30\textwidth}*{5}{R}@{}}
\toprule
& € Euro equities / 欧元欧洲股票 & \$ US equities / 美元美国股票 & £ UK equities / 英镑英国股票 & Japan equities / 日本股票 & \$ EM equities / 美元新兴市场股票 \\
\midrule
Daily volatility target / 日波动率目标 (A) & €42,780.47 & €42,780.47 & €42,780.47 & €42,780.47 & €42,780.47 \\
Price / 价格 (B) & €27 & \$30 & £5 & €13 & \$38 \\
Number of shares per block / 每单位股数 (C) & 100 & 100 & 100 & 100 & 100 \\
FX / 汇率 (D) & 1 & 0.9 & 1.2 & 1 & 0.9 \\
Block value / 单位价值 (E) & €27 & \$30 & £5 & €13 & \$38 \\
\emph{C × B ÷ 100} & & & & & \\
Price volatility / 价格波动率, \% (G) & 1.90 & 2.00 & 2.00 & 1.70 & 2.50 \\
Instrument currency volatility / 品种货币波动率 (H) & €51.3 & \$60 & £10 & €22.1 & \$95 \\
\emph{G × E} & & & & & \\
Instrument value volatility / 品种价值波动率 (I) & €51.3 & €54 & €12 & €22.1 & €85.5 \\
\emph{H × D} & & & & & \\
Volatility scalar / 波动率缩放因子 (J) & 834 & 792 & 3565 & 1936 & 500 \\
\emph{A ÷ I} & & & & & \\
\bottomrule
\end{tabularx}
}
\end{chinesetableblock}

\bilingualsubsection{Equity Calculations Part Two}{股票部分计算：第二部分（再次更新）}

\begin{chinesetableblock} \noindent{\scriptsize \setlength{\tabcolsep}{2pt} \renewcommand{\arraystretch}{1.12} \begin{tabularx}{\textwidth}{@{} p{0.30\textwidth}*{5}{R}@{}}
\toprule
& € Euro equities / 欧元欧洲股票 & \$ US equities / 美元美国股票 & £ UK equities / 英镑英国股票 & Japan equities / 日本股票 & \$ EM equities / 美元新兴市场股票 \\
\midrule
Forecast / 预测值 (K) & +10 & +10 & +10 & +10 & +10 \\
Subsystem position, blocks / 子系统仓位，单位 (L) & 834 & 792 & 3565 & 1936 & 500 \\
\emph{K × J ÷ 10} & & & & & \\
Instrument weight / 品种权重 (M) & 9.02\% & 9.02\% & 9.02\% & 9.02\% & 15.47\% \\
Instrument diversification multiplier / 品种分散化乘数 (N) & 1.61 & 1.61 & 1.61 & 1.61 & 1.61 \\
Portfolio instrument position, blocks / 投资组合品种仓位，单位 (O) & 121.1 & 115.1 & 517.8 & 281.2 & 124.6 \\
\emph{L × M × N} & & & & & \\
Rounded target position, blocks / 四舍五入后的目标仓位，单位 (P = round O) & 121 & 115 & 518 & 281 & 125 \\
Current position, blocks / 当前仓位，单位 & 195 & 200 & 786 & 528 & 200 \\
Trade, blocks / 交易，单位 & Sell 74 / 卖出 74 & Sell 85 / 卖出 85 & Sell 268 / 卖出 268 & Sell 247 / 卖出 247 & Sell 75 / 卖出 75 \\
\bottomrule
\end{tabularx}
}
\end{chinesetableblock}

\begin{englishblock} Notice how most of the equity selling is caused by increases in price volatility, rather than the changes to instrument weights. That draws this example to an end as you've now seen most of the significant trading action, although it's tempting to continue and let the bond overweight come good. \end{englishblock}

\begin{chineseblock}请注意，这一轮股票卖出的主要原因，其实是价格波动率上升，而不是品种权重变化本身。这个例子就到此为止吧，因为你已经看到了大部分最关键的交易动作，尽管从故事角度上看，继续持有债券超配直到它发挥作用也同样很诱人。\end{chineseblock}
